\documentclass{article}

\usepackage[width=14cm, left=3cm, top=2cm]{geometry}

\usepackage[brazil]{babel}
\usepackage[T1]{fontenc}
\usepackage[utf8]{inputenc}

%\usepackage{fontspec}
%\setmainfont{Noto Serif}

\usepackage{amsmath}
\usepackage{amssymb}

\usepackage{minted}

\usepackage{hyperref}
\hypersetup{
    colorlinks=true,
    breaklinks=true,
    linkcolor=purple,
    filecolor=cyan,
    urlcolor=blue,
    pdftitle={Picross e Ant Colony},
}
\urlstyle{same}

\usepackage{subcaption}
\usepackage{graphicx}
\graphicspath{{../images/}}

\newcommand{\apl}{\texttt{APL}}
\newcommand{\lang}{\texttt{J}}
\newcommand{\jver}{\texttt{J903}}
\newcommand{\inlcode}{\mintinline[breaklines]{J}}

\newcommand{\file}{\texttt}

\title{Trabalho Computacional - Programação Não Linear\\
Picross e Ant Colony}
\author{Daniel Kiyoshi Hashimoto Vouzella de Andrade - 119025937}
\date{Novembro 2022}

\begin{document}
\maketitle

\section{Problema: Picross}

O problema escolhido foi a resolução de
um puzzle de \emph{picross}.
Também conhecido por vários outros nomes
(nonogram, griddler, hanjie, pic-a-pix,
paint by numbers, \dots),
o puzzle tem como objetivo colorir
uma grid \(n \times m\) com quadrados,
seguindo \(n + m\) dicas
(exatamente uma para cada linha e coluna).
Por simplificação, só vamos considerar
puzzles de tamanho \(n \times n\).

Cada dica é uma lista de números
que indicam o tamanho dos blocos e
sua ordem da sua respectiva linha ou coluna,
e assim o tamanho da lista
indica a quantidade de blocos dessa linha/coluna.

A figura \ref{pg:5x5} mostra dois exemplos
de puzzles \(n \times n\) já resolvidos,
mas note que as dicas não estão incluídas e
a borda cinza não faz parte do puzzle.
As dicas são representadas por uma matriz
\(2 \times n\) de lista de números
(eles estão ``encaixotados''),
de forma que a lista \((0,i)\)
indica a ordem e tamanho dos blocos
(esquerda para direita)
da linha \(i\)
(contando de cima para baixo).
E, de forma similar,
a lista \((1,i)\)
indica a ordem e tamanho dos blocos
(cima para baixo)
da coluna \(i\)
(contando da esquerda para direita).
Em ambos os casos \(0 \le i < n\).
Caso não tenha nenhum bloco na linha/coluna,
aparece um \(0\) no lugar de uma lista vazia.

A seguir, estão as dicas dos puzzles
\ref{p:face}, \ref{p:heart} e \ref{p:random5};
a linha que está identada é o código em \lang{}
e abaixo o resultado:
\begin{minted}{j}
    NB. $ é usado para alterar as dimensões do array
    2 5 $ 1 ; 1 1 ; 1 ; 1 1 ; 1 ; 0 ; 1 1 ; 0 ; 1 1 ; 3
+-+---+-+---+-+
|1|1 1|1|1 1|1|
+-+---+-+---+-+
|0|1 1|0|1 1|3|
+-+---+-+---+-+
   2 5 $ 1 1 ; 5 ; 5 ; 3 ; 1 ; 2 ; 4 ; 4 ; 4 ; 2
+---+-+-+-+-+
|1 1|5|5|3|1|
+---+-+-+-+-+
|2  |4|4|4|2|
+---+-+-+-+-+
   2 5 $ 2 2 ; 1 1 ; 1 1 ; 1 1 ; 1 ; 2 1 ; 1 2 ; 1 ; 1 2 ; 1
+---+---+---+---+-+
|2 2|1 1|1 1|1 1|1|
+---+---+---+---+-+
|2 1|1 2|1  |1 2|1|
+---+---+---+---+-+
\end{minted}

\begin{figure}[h]
    \centering
    \begin{subfigure}{0.3\textwidth}
        \centering
        \includegraphics[width=0.9\linewidth]{face-border}
        \caption{Carinha Feliz de ASCII}
        \label{p:face}
    \end{subfigure}
    \begin{subfigure}{0.3\textwidth}
        \centering
        \includegraphics[width=0.9\linewidth]{heart-border}
        \caption{Coração}
        \label{p:heart}
    \end{subfigure}
    \begin{subfigure}{0.3\textwidth}
        \centering
        \includegraphics[width=0.9\linewidth]{random5-border}
        \caption{Aleatório}
        \label{p:random5}
    \end{subfigure}
    \caption{Exemplos de picross 5x5}
    \label{pg:5x5}
\end{figure}

Ambos os problemas de existência e unicidade de uma solução
(saber se existe uma solução e se ela é única)
de um puzzle arbitrário de picross são \emph{NP-completo}.

\subsection{Formas conhecidas de resolução}
\begin{itemize}
    \item Depth First Search (Brute Force): \par
        Tenta pintar um pixel,
        se nenhuma dica quebrar tente de novo.
        Se alguma dica quebrar, troque a cor e continue.
        Se der errado com a nova cor, ``volte'' e tente de novo.
    \item Heuristic/Iterational: \par
        A ideia é aplicar uma sequência de regras simples
        para determinar a cor de algum pixel
        várias e várias vezes.
        A técnica é semelhante a um humano resolvendo o puzzle.
        Alguns puzzles mais complexos
        não podem ser resolvidos dessa forma,
        é comum usar um DFS quando
        uma iteração não encontra nenhum pixel novo.
    \item Integer Linear Programming: \par
        O puzzle é modelado como
        uma sequência de restrições lineares
        com uma função \emph{dummy} de otimização.
\end{itemize}

\section{Modelagem}

O algoritmo escolhido foi o \emph{Ant Colony}.
O principal motivo para ele ter sido escolhido
foi ``formigas são legais!'',
na verdade, ele foi escolhido antes do problema em si.
O problema é muito bom para essa heurística,
pois é o problema é facilmente traduzido
em escolher um caminho sobre um grafo e
grafo é gerado bem simples.
\emph{Ant Colony} funciona com um simples loop,
no qual a cada iteração são executados os passos:
\begin{enumerate}
    \item Construir caminhos \par
        Nesse passo, são construídos \(ants\) caminhos e
        o melhor deles, de acordo com a \emph{função objetivo},
        é guardado.
        Cada caminho, deve ser escolhido aleatóriamente
        em função da \emph{distribuição do feromônio} e
        da \emph{heurística}.
        Geralmente são usados os parâmetros \(\alpha\) e \(\beta\)
        para equilibrar esses valores
        que formam os pesos dos caminhos.
    \item Atualizar o feromônio \par
        O melhor caminho da iteração \(c\) é usado
        para atualizar
        o feromônio \(\tau_e\) em cada aresta \(e\) do grafo
        usando a seguinte fórmula:
        \[
            \tau_e \leftarrow
                (1 - \rho) \; \tau_e
                + in(c, e) \; \frac{\rho \; \tau}{len(c)}
        \]
        onde \(\rho\) é a taxa de evaporação do feromônio
        (um parâmetro),
        \(\tau\) é a soma de feromônio em todas as atestas
        (uma constante),
        \(in(c, e)\) é uma função que retorna \(1\) se
        \(c\) passa pela aresta \(e\) e \(0\) caso contrário e
        \(len(c)\) são por quantas arestas \(c\) passa.
    \item Checar a condição de parada \par
        Existem duas condições mais comuns:
        checar se ``acabou o gás'',
        podendo ser o número de iterações ou tempo de execução;
        checar se achamos uma solução.
        Nem sempre é possível ou fácil usar a segunda,
        mas nesse caso é possível e relativamente fácil.
        Quando for decidido que se deve parar,
        retorna-se o melhor caminho encontrado até agora
        (não necessariamente dessa iteração).
\end{enumerate}

\subsection{Função Objetivo}
A \emph{função objetivo} deve impor
uma ordenação sobre os caminhos,
de forma que quanto mais próxima da solução
``melhor'' ela é.

\subsection{Heurística}
A \emph{heurística} é uma forma de aumentar a chance de escolher
um caminho mais provável de ser melhor.
Também deixa pouco provável
caminhos que aparentam ser ruins.
Idealmente ela vai apontar para caminhos que
minimizam/maximizam a \emph{função objetivo}.

\section{Implementação}

A linguagem escolhida para a implementação foi \lang{}:
uma linguagem orientada a arrays,
da mesma família que \apl{}.
Não foi usada nenhuma biblioteca para construir o algoritmo,
mas foi usada \inlcode{viewmat} para fazer as imagens
apresentadas nesse relatório.

\subsection{Representação do Grafo e das Dicas}
Para um puzzle \(n \times m\),
construímos um grafo com \(1 + n \cdot m\) vértices
numerados de \(-1 \le v < n \cdot m\) e
\(2 \cdot n \cdot m\) arestas nomeadas \((p, b)\),
com \(0 \le p < n \cdot m\) e \(b \in \{0, 1\}\).

\(-1\) é o vértice inicial do caminho,
só está alí por formalidades;
os outros \(0 \le v < n \cdot m\)
representam um pixel de coordenadas \((x(v), y(v))\) do puzzle.
A princípio, \(x\) e \(y\) podem ser quaisquer
funções que, em conjunto,
não repetem nem deixam sobrar nenhum pixel.
Por simplicidade, vamos usar
\(x(v) := \!\!\!\mod\!\!(v, m)\) e
\(y(v) := \left\lfloor\frac{v}{m}\right\rfloor\),
mas, pela natureza da linguagem escolhida,
isso vai ``aparecer'' como a ausência de uma função
que reordena uma possível solução no código.

A aresta \((p, b)\) sai do vértice \(p-1\) e
chega no vértice \(p\),
\(b\) indica a cor do pixel \(p\):
se \(b = 0\) é vazio/branco,
é preto caso contrário.
Na prática podemos ``jogar os vértices fora''
e só prestar atenção nas arestas,
isso torna grafo é tão simples
a ponto de não aparecer explicitamente na implementação.

A representação das dicas já foi descrita na primeira seção,
mas não foi dito que ela só funciona para puzzles quadrados
(todas as linhas de uma matriz devem ter o mesmo tamanho).
Não é difícil criar uma outra representação que suporte
puzzles retangulares, basta encaixotar cada lista,
mas não tinha pensado em suportar puzzles retangulares durante
a implementação.
A dica do puzzle \ref{p:face} ficaria assim:
\begin{minted}{j}
   (1 ; 1 1 ; 1 ; 1 1 ; 1) ; < 0 ; 1 1 ; 0 ; 1 1 ; 3
+---------------+---------------+
|+-+---+-+---+-+|+-+---+-+---+-+|
||1|1 1|1|1 1|1|||0|1 1|0|1 1|3||
|+-+---+-+---+-+|+-+---+-+---+-+|
+---------------+---------------+
\end{minted}

\subsection{Representação de um Caminho}
O caminho é representado por uma matriz de mesma forma
do puzzle \(n \times m\) com os valores em \(\{ 0, 1 \}\).
O elemento na posição \((i, j)\) da matriz que guarda o valor \(b\)
representa a aresta \((xy^{-1}(i, j), b)\),
onde \(xy^{-1}\) é a ``função inversa'' das funções \(x\) e \(y\),
ou seja, \(xy^{-1}(i, j) = v\) quando \(x(v) = i\) e \(y(v) = j\).
No nosso caso, \(xy^{-1}(i, j) := i \cdot m + j\).

Essa representação é muito interessante, pois
``a menos de uma função de reorganização''
(no nosso caso, a ``reorganização vazia'')
o caminho é uma possível representação direta da solução.
Veja a representação da solução do puzzle \ref{p:heart}:
\begin{minted}{j}
    NB. Isso é só para guardar numa variável e mostrar o resultado
    ] b =: 0 1 0 1 0 , 1 1 1 1 1 , 1 1 1 1 1 , 0 1 1 1 0 ,: 0 0 1 0 0
0 1 0 1 0
1 1 1 1 1
1 1 1 1 1
0 1 1 1 0
0 0 1 0 0
    NB. 'X' é preto e '.' é branco
    b { '.X'
.X.X.
XXXXX
XXXXX
.XXX.
..X..
\end{minted}

\subsection{Função Objetivo}
A \emph{função objetivo} escolhida transforma
o caminho em uma matriz de lista de tamanho dos blocos,
de forma que,
se essa fosse a solução retornaria uma matriz idêntica a dica.
Então, as listas são pareadas e elemento a elemento
(possívelmente estendendo a lista menor com \(0\)s)
são feitas subtrações seguidas de absoluto (para remover negativos).
Somamos todos os números de todas as listas.
A função que faz isso é \inlcode{flist}.

Um exemplo comparando as segundas linhas dos
puzzles \ref{p:heart} e \ref{p:random5}:
\begin{minted}{j}
    5 ; 1 1         NB. em representação de blocos
+-+---+
|5|1 1|
+-+---+
    5 0 ,:!.0 1 1   NB. estendendo a menor com 0s
5 0
1 1
    | 5 0 - 1 1     NB. subtraindo, seguido de absoluto
4 1
    +/ 4 1          NB. somando (falta as outras linhas e colunas)
5
\end{minted}

Com essa escolha de função,
procuramos minimizar a sua imagem.
Uma propriedade interessante dela é
que apenas uma solução vai ter o valor de \(0\),
permitindo parar mais cedo.

\subsection{Representação do Feromônio}
A lista de feromônios é representada por uma
matriz \(2 \times n \times m\),
no qual o elemento \((b, i, j)\) armazena a
quantidade de feromônio da aresta \((xy^{-1}(i, j), b)\),
usando o \(xy^{-1}\) descrito anteriormente.

Colocar o \(b\) na primeira dimensão
é vantajoso pela semântica da linguagem.
Essa representação permite pegar
uma matriz que representa um caminho \(c\)
e ``colar em cima'' sua matriz com elementos complementares
(\(x \longmapsto 1 - x\)),
montando as respostas da função \(in(c, e)\)
(uma ``máscara'')
necessárias para atualizar o feromônio.
Também facilita na geração de caminhos,
explicada mais a frente.

\subsection{Construção de Caminhos}
Não foi escolhida nenhuma heurística
para a construção dos caminhos,
o principal motivo é simplificar a implementação
dado que estou fazendo o trabalho sozinho e
não ter pensado em uma heurística fácil de implementar.

Dessa forma, não precisamos
dos parâmetros \(\alpha\) e \(\beta\),
e a escolha do caminho,
pelo método da roleta,
é influenciada apenas pelo feromônio.
Isso, junto a escolha do grafo,
faz com que a escolha da cor de um pixel
seja completamente independente
da escolha da cor de qualquer outro pixel.
O esperado é que seja necessário mais iterações
para que se ache um caminho bom.

\subsection{Condição de Parada}
A condição de parada é a combinação
de um contador de iterações
com uma verificação da imagem da \emph{função objetivo},
se algum deles for \(0\) podemos parar.

\lang{} possui uma função de ``exponenciação de funções'',
que aplica a função a quantidade de vezes que o ``expoente'' indica,
de forma que para um \(n\) inteiro:
\[
    f^n(x) := \begin{cases}
        f^{n+1}(f^{-1}(x))  &\quad,\; n < 0 \\
        x                   &\quad,\; n = 0 \\
        f^{n-1}(f(x))       &\quad,\; n > 0
    \end{cases}
\]
Não precisa se preocupar com o inverso da função.
O interessante é que está definida para \(n = \infty\),
aplica-se \(f\) até o resultado não se alterar,
efetivamente calculando um \emph{ponto fixo} da função.
Na implementação, o loop principal retorna
a mesma entrada caso uma das condições indique para parar,
permitindo usar essa ideia de \emph{ponto fixo}.

O número máximo de iterações poderia ser \(n\) (o expoente),
no lugar de \(\infty\),
mas deixá-lo como um parâmetro
permite uma \inlcode{main} mais simples.

\section{Resultados}

Na seção \ref{s:bib},
tem um link com instruções de instalação.
A versão usada nesse trabalho foi a \jver{},
mas não deve ter grandes problemas se a versão for diferente.

Vamos seguir assumindo que já estamos em um sessão
(\texttt{jconsole}, \texttt{jqt} e \texttt{jhs} são exemplos),
e guardamos o caminho
até o \file{main.ijs} na variável \inlcode{path}.
Então, chamamos o verbo \inlcode{load}
para carregar as funções que estão no arquivo \file{main.ijs}.
Por exemplo:
\begin{minted}{j}
    NB. ']' é só para mostrar o valor da variável
    ] path =: '~/path/to/main.ijs'
~/path/to/main.ijs
    load path
    NB. Note que: load '~/path/to/main.ijs'
    NB. também funciona
\end{minted}

\subsection{Como rodar}
A função \inlcode{main} recebe
uma lista encaixotada de \emph{parâmetros},
vamos guardar em \inlcode{p}, e
uma lista encaixotada de \emph{variáveis de loop},
vamos guardar em \inlcode{v}.

Os \emph{parâmetros} são, em ordem:
\emph{taxa de evaporação} \(\rho\),
\emph{função objetivo} \(f\) e
\emph{número de formigas por iteração} \(ants\).
O exemplo assume que \(\rho = 0.4\),
\(f\) é está otimizando para o puzzle \ref{p:face} e
\(ants = 100\).
\begin{minted}{j}
    face
0 0 0 1 0
0 1 0 0 1
0 0 0 0 1
0 1 0 0 1
0 0 0 1 0
    n2b face    NB. calcula as dicas a partir de uma solução
+-+---+-+---+-+
|1|1 1|1|1 1|1|
+-+---+-+---+-+
|0|1 1|0|1 1|3|
+-+---+-+---+-+
    ] p =: (<0.4)`((n2b face)&f)`(<100)
+---+---------------------------+---+
|0.4|+-+-----------------------+|100|
|   ||&|+-------------------+-+||   |
|   || ||+-+---------------+|f|||   |
|   || |||0|+-+---+-+---+-+|| |||   |
|   || ||| ||1|1 1|1|1 1|1||| |||   |
|   || ||| |+-+---+-+---+-+|| |||   |
|   || ||| ||0|1 1|0|1 1|3||| |||   |
|   || ||| |+-+---+-+---+-+|| |||   |
|   || ||+-+---------------+| |||   |
|   || |+-------------------+-+||   |
|   |+-+-----------------------+|   |
+---+---------------------------+---+
    NB. Essa caixa do meio é a representação da função objetivo
\end{minted}

As \emph{váriavei de loop} são, em ordem:
\emph{contador de iterações} \(n\)
que conta quantas iterações ainda pode seguir,
\emph{lista de feromônio} \(\tau\),
\emph{imagem do caminho \emph{dummy}} e
\emph{caminho \emph{dummy}},
um caminho todo em branco para
servir como um \emph{``placeholder''}.
Elas podem ser criadas
com uma função auxiliar \inlcode{init}.
O exemplo assume que \(n = 200\) e
que a \(\tau\) começa com \(1\) em cada aresta.
\begin{minted}{j}
    NB. não precisa ser a solução,
    NB. basta uma matriz com o tamanho igual
    ] fer =: 1 newfer face
1 1 1 1 1
1 1 1 1 1
1 1 1 1 1
1 1 1 1 1
1 1 1 1 1

1 1 1 1 1
1 1 1 1 1
1 1 1 1 1
1 1 1 1 1
1 1 1 1 1
   ] v =: 200 (p init) fer
+---+---------+--+---------+
|200|1 1 1 1 1|14|0 0 0 0 0|
|   |1 1 1 1 1|  |0 0 0 0 0|
|   |1 1 1 1 1|  |0 0 0 0 0|
|   |1 1 1 1 1|  |0 0 0 0 0|
|   |1 1 1 1 1|  |0 0 0 0 0|
|   |         |  |         |
|   |1 1 1 1 1|  |         |
|   |1 1 1 1 1|  |         |
|   |1 1 1 1 1|  |         |
|   |1 1 1 1 1|  |         |
|   |1 1 1 1 1|  |         |
+---+---------+--+---------+
\end{minted}

O algoritmo propriamente dito,
está implementado em \inlcode{loop},
mas a \inlcode{main} coleta mais informações e
é mais fácil de usar.

A função \inlcode{main} pode ser chamada com os argumentos
\inlcode{p} pela esquerda e \inlcode{v} pela direita.
Antes de chamá-la, pode-se alterar a semente inicial
se achar interessante.
\inlcode{main} retorna uma lista de caixas,
a primeira caixa contém a semente inicial
com o estado do gerador de números aleatórios
antes de rodar o algoritmo,
mas esse estado é gigantesco,
então não vai ser mostrado no exemplo;
a segunda o tempo de execução;
a terceira é uma variável interna que indica se deve parar,
já que ele retornou algo, deve ser sempre \(0\);
a quarta são as \emph{variáveis de loop}:
\(n\), \(\tau\), \(f(c^*)\), \(c^*\).
\begin{minted}{j}
    NB. Escolhendo uma semente (opcional)
    wseed 16807

    'rng time stop ret' =: p main v
    time ; stop ; <ret
+-------+-+------------------------------------------------------------+
|2.03215|0|+-+--------------------------------------------+-+---------+|
|       | ||0| 1.32328 0.929616 0.620616   1.4589 0.859697|4|1 1 0 1 0||
|       | || |0.784592  1.43173  1.20543  1.51852  1.61536| |1 1 1 1 1||
|       | || | 1.36438  1.78557 0.427282  1.19983  1.08449| |1 1 0 0 0||
|       | || |0.632997 0.864541 0.512813 0.708474  1.00561| |0 0 0 1 0||
|       | || |0.453377 0.409504  1.34556 0.847346  1.18167| |0 1 1 1 0||
|       | || |                                            | |         ||
|       | || |0.676724  1.07038  1.37938 0.541097   1.1403| |         ||
|       | || | 1.21541 0.568269 0.794566 0.481477 0.384644| |         ||
|       | || |0.635621 0.214428  1.57272 0.800168 0.915508| |         ||
|       | || |   1.367  1.13546  1.48719  1.29153 0.994391| |         ||
|       | || | 1.54662   1.5905 0.654442  1.15265 0.818328| |         ||
|       | |+-+--------------------------------------------+-+---------+|
+-------+-+------------------------------------------------------------+
    NB. Se não quiser guardar a saída em variáveis pode fazer:
    NB. }. p main v
    NB. para jogar fora o 'rng', ou
    NB. 1 3 { p main v
    NB. para escolher quais resultados quer ver
    NB. (nesse caso 'time' e 'ret')
\end{minted}

\subsection{Método}

\section{Possíveis Melhorias de Performance}
\subsection{Representação do Grafo}

\subsection{Função Objetivo}
\subsubsection{Quadrado em vez de Absoluto}
\subsubsection{Função mais inteligente}

\subsection{Geração de Caminhos}

\section{Possíveis Melhorias de Funcionalidade}
\subsection{Permitir para diferentes linhas e colunas}
\subsection{Generalizar para mais cores}
\subsection{Permitir solução parcial}

\section{Análise final}

\section{Bibliografia} \label{s:bib}

Sobre \lang{}:
\begin{itemize}
    \item Página inicial: \par
        \url{https://wiki.jsoftware.com/wiki/Main_Page}
    \item Instruções de instalação: \par
        \url{https://wiki.jsoftware.com/wiki/System/Installation}
    \item Dicionário de primitivas: \par
        \url{https://wiki.jsoftware.com/wiki/NuVoc}
    \item J-playground (uma sessão online): \par
        \url{https://jsoftware.github.io/j-playground/bin/html/index.html}
\end{itemize}

Sobre \emph{Picross}:
\begin{itemize}
    \item History of Picross (Nonograms): \par
        \url{https://youtu.be/qJCPxyi5x5g}
    \item Solving Colored Nonograms: \par
        \url{https://run.unl.pt/bitstream/10362/2388/1/Mingote_2009.pdf}
    \item NP Completeness of Nonograms: \par
        \url{https://kelbybsandvick.github.io/pdf/NPProblemPaper.pdf}
    \item Complexity and solvability of Nonogram puzzles: \par
        \url{https://fse.studenttheses.ub.rug.nl/15287/1/Master_Educatie_2017_RAOosterman.pdf}
    \item kamilkhanlab/nonogram-ilp
        (Implementação de solver em GAMS): \par
        \url{https://github.com/kamilkhanlab/nonogram-ilp}
    \item Nonogram Wikipédia: \par
        \url{https://en.wikipedia.org/wiki/Nonogram}
\end{itemize}

\end{document}
